\chapter{Riesgos y plan de ensayos}\label{cap:riesgos}

\section{Análisis de modos de falla (AMFE)}
Severidad (S), ocurrencia (O) y detección (D) de 1 a 10; el número de prioridad es
$\mathrm{NPR}=S\cdot O\cdot D$. Los valores son estimaciones del equipo antes de los
ensayos y se actualizan con cada resultado.

\begin{longtable}{@{}p{3.3cm}p{2.7cm}ccccp{4.6cm}@{}}
\caption{AMFE del sistema de descenso, ordenado por NPR.}\label{tab:amfe}\\
\toprule
\textbf{Modo de falla} & \textbf{Efecto} & \textbf{S} & \textbf{O} & \textbf{D} & \textbf{NPR} & \textbf{Mitigación}\\\midrule
\endfirsthead
\toprule
\textbf{Modo de falla} & \textbf{Efecto} & \textbf{S} & \textbf{O} & \textbf{D} & \textbf{NPR} & \textbf{Mitigación}\\\midrule
\endhead
\bottomrule
\endfoot
El rotor no arranca & Caída a $\approx\SI{13}{m/s}$; falla \Req{C4} & 9 & 4 & 4 & 144 & Paso $\le\ang{-2}$ (tabla~\ref{tab:arranque}); ensayo E2 de 10 arranques; reintento del servo\\
La traba no suelta & Rotor no se despliega & 9 & 3 & 3 & 81 & Servo con engranajes metálicos; 20 liberaciones en E1; reintento automático\\
Drogue en la estela del rotor & Caída hasta \SI{7.8}{m/s} & 5 & 5 & 3 & 75 & Cuerda de \SI{1.5}{m}; E4 mide con y sin drogue\\
Liberación prematura & Etapa 1 corta; \Req{C4} dudoso & 7 & 2 & 5 & 70 & Solo en \texttt{DESCENT}; $h_{\max}$ en FRAM; filtro de 50\,Hz\\
Palas que se abren en el ascenso & Daño, despliegue fuera de tiempo & 8 & 2 & 4 & 64 & Traba positiva; espuma detrás de las palas; E5 vibración\\
Cuerda que se retuerce o roza & Rotura de la cuerda & 7 & 2 & 4 & 56 & Destorcedor; borde del eje redondeado y pulido\\
Desbalance de palas & Vibración, video inutilizable & 4 & 4 & 2 & 32 & Balanceo a $\pm\SI{0.1}{g}$; palas marcadas\\
Conicidad excesiva o flameo & Menos área efectiva, rotura & 6 & 2 & 3 & 36 & Tope a \ang{+10}; centro de masa al 25\,\%\\
Masa fuera de \SI{500 \pm 10}{g} & Falla \Req{S1} & 6 & 3 & 1 & 18 & Presupuesto en CAD; lastre ajustable\\
\end{longtable}

\section{Ensayos}

\begin{table}[H]
\centering\small
\caption{Plan de ensayos, en orden de ejecución.}\label{tab:ensayos}
\begin{tabularx}{\textwidth}{@{}clYY@{}}
\toprule
\textbf{N.º} & \textbf{Ensayo} & \textbf{Qué se mide} & \textbf{Criterio de aceptación}\\\midrule
E1 & Banco: mecanismo & Giro libre a mano, balanceo, 20 liberaciones de la traba & Gira libre $>\SI{5}{s}$; 20 de 20 liberaciones\\
E2 & Ventilador vertical (o auto a \SI{20}{km/h}) & Arranque desde cero, rpm, conicidad (video), paso óptimo & 10 de 10 arranques en $<\SI{5}{s}$; rpm medidas $\pm$20\,\% de 1150\\
E3 & Balanza de empuje en túnel o ventilador & Empuje frente a velocidad del aire: obtiene $\CR$ real & $\CR\ge1{,}0$\\
E4 & Caída desde drone (\SIrange{30}{60}{m}) & Velocidad con barómetro, con y sin drogue & $\Vr\le\SI{7.5}{m/s}$ sin drogue; $\le\SI{7}{m/s}$ con drogue\\
E5 & Vibración y choque & 15\,G y 30\,G (\Req{S4}, \Req{S5}) & Nada se suelta; traba sigue cerrada\\
E6 & Integración completa & Secuencia de estados, video, telemetría & Todos los estados en el CSV; video de las dos etapas\\
\bottomrule
\end{tabularx}
\end{table}

\begin{figure}[H]
\centering
\begin{tikzpicture}[x=1mm, y=1mm]
  % ---- E2: ventilador vertical
  \node[font=\small\bfseries] at (35,92) {E2: ventilador vertical};
  \draw[fijo] (10,0) rectangle (60,12);
  \node[font=\scriptsize] at (35,6) {ventilador};
  \foreach \x in {18,26,34,42,50}{\draw[flecha, cfijo!60] (\x,14) -- (\x,32);}
  \draw[cgris] (35,12) -- (35,40);
  \draw[black, line width=0.8pt] (33,40) -- (37,40);
  \filldraw[fijo] (31,40) rectangle (39,62);
  \filldraw[rot] (14,63) rectangle (56,64.5);
  \filldraw[rot] (30,62) rectangle (40,66);
  \draw[cgris] (35,66) -- (35,84);
  \draw[black] (28,84) -- (42,84);
  \node[font=\scriptsize, anchor=west] at (43,84) {celda de carga};
  \node[font=\scriptsize, anchor=west] at (58,63) {rotor};
  \node[font=\scriptsize, anchor=west] at (40,28) {anemómetro};
  \draw[black] (40,24) circle (1.2);
  \node[font=\scriptsize, anchor=west] at (40,50) {tacómetro óptico};
  \draw[crot] (14,63) -- (6,70);
  % ---- E4: caída desde drone
  \begin{scope}[shift={(85,0)}]
  \node[font=\small\bfseries] at (35,92) {E4: caída desde drone};
  \draw[fijo] (22,80) rectangle (48,84);
  \draw[fijo] (18,84) -- (52,84);
  \node[font=\scriptsize] at (35,87) {drone};
  \draw[cgris, densely dashed] (35,80) -- (35,74);
  \draw[act] (33,72) rectangle (37,74);
  \node[font=\scriptsize, anchor=west] at (38,73) {liberador};
  \filldraw[fijo] (32,45) rectangle (38,60);
  \filldraw[rot] (20,60) rectangle (50,61.2);
  \draw[cuerda] (35,61) -- (35,70);
  \draw[flecha, cgris] (60,60) -- (60,10) node[midway, right, font=\scriptsize, align=left] {30--60 m\\ barómetro\\ a 50\,Hz};
  \draw[black] (0,5) -- (70,5);
  \node[font=\scriptsize] at (35,2) {suelo};
  \end{scope}
\end{tikzpicture}
\caption{Esquemas de los ensayos E2 y E4.}
\label{fig:ensayos}
\end{figure}

\section{Cronograma}
\begin{figure}[H]
\centering
\begin{ganttchart}[hgrid, vgrid, x unit=0.82cm, y unit chart=0.62cm,
    title label font=\scriptsize, bar label font=\small,
    bar/.append style={fill=cfijo!35, draw=cfijo},
    milestone/.append style={fill=cfreno, draw=cfreno}]{1}{14}
  \gantttitlelist{1,...,14}{1}\\
  \ganttbar{CAD y presupuesto de masa}{1}{2}\\
  \ganttbar{Fabricación de palas y cubo}{2}{4}\\
  \ganttbar{E1 banco}{4}{4}\\
  \ganttbar{E2 ventilador, elección de paso}{5}{6}\\
  \ganttbar{E3 empuje ($\CR$ real)}{6}{7}\\
  \ganttmilestone{Rotor congelado}{7}\\
  \ganttbar{Electrónica y software de estados}{3}{8}\\
  \ganttbar{E4 caídas desde drone}{8}{9}\\
  \ganttbar{E5 vibración y choque}{10}{10}\\
  \ganttbar{E6 integración}{11}{12}\\
  \ganttmilestone{Listo para vuelo}{12}\\
  \ganttbar{Margen}{13}{14}
\end{ganttchart}
\caption{Cronograma tentativo en semanas desde el inicio del trabajo.}
\label{fig:gantt}
\end{figure}
